\section{KVM v2 Design}

\subsection{Model}

\KVMvTwo{} follows the gVisor KVM-platform idea, adapted to \UML{}:
guest userspace runs at ring 3 inside a KVM guest, the \UML{} kernel
runs at ring 0 inside the same KVM guest, and ordinary syscalls enter
the \UML{} kernel through MSR\_LSTAR instead of trapping to a host
monitor \cite{gvisor-platforms,kvm-api}. Host exits are still used for
operations that truly require host participation, but the target is
to keep the common syscall path in guest context.

\begin{figure}[h]
\centering
\begin{tikzpicture}[
  region/.style={draw=UMLGrid,rounded corners=3pt,align=center,inner sep=8pt},
  arr/.style={-{Latex[length=2mm]},thick,UMLMuted},
]
\node[region,fill=UMLGreen!10,minimum width=3.8cm,minimum height=1.25cm] (user) {guest userspace\\ring 3};
\node[region,fill=UMLBlue!10,minimum width=3.8cm,minimum height=1.25cm,right=1.1cm of user] (kernel) {\UML{} kernel\\ring 0};
\node[region,fill=UMLPanel,minimum width=8.7cm,minimum height=1.0cm,below=0.7cm of $(user)!0.5!(kernel)$] (host) {host process, KVM ioctls, signals, files, timers};
\draw[arr] (user) -- node[above,font=\footnotesize]{\texttt{syscall}/LSTAR} (kernel);
\draw[arr] (kernel) -- node[right,font=\footnotesize]{VMCALL, HLT, MMIO, host work} (host);
\draw[arr] (host) -| node[pos=0.25,below,font=\footnotesize]{resume} (user);
\end{tikzpicture}
\caption{\KVMvTwo{} execution model. It is not a QEMU-style machine; it is a trap mechanism for \UML{} itself.}
\end{figure}

\subsection{Implementation Themes}

The current \KVMvTwo{} design is organized around five hard problems:

\begin{enumerate}
  \item \textbf{Guest entry and return.} The backend must enter guest
        ring 3, handle SYSCALL/SYSRET transitions, and recover when
        host signals interrupt critical regions.
  \item \textbf{Per-vCPU and per-task state.} General registers,
        CR2, segment bases, FPU/XSAVE, pending events, IST frames,
        and KVM run state cannot leak across tasks or host CPUs.
  \item \textbf{Memory and TLB state.} \UML{} page-table changes must
        be reflected into KVM-visible mappings and cross-vCPU TLB
        state without stale translations.
  \item \textbf{SMP correctness.} Host scheduling, guest scheduling,
        signal delivery, and KVM vCPU ownership must compose without
        assuming that \texttt{preempt\_disable()} pins a host thread.
  \item \textbf{Performance.} The backend must preserve the win from
        in-guest dispatch while avoiding correctness regressions from
        lazy state handling.
\end{enumerate}

\subsection{Design Detail: Main Control Surfaces}

\begin{table}[h]
\centering
\caption{\KVMvTwo{} implementation surfaces to understand before reviewing patches.}
\small
\begin{tabularx}{\linewidth}{@{}L{0.21\linewidth}L{0.29\linewidth}Y@{}}
\toprule
Surface & Representative code & Design role \\
\midrule
Initialization and capabilities & \path{init.c}, \path{ops.c}, \path{Kconfig} & Probe host KVM support, wire backend ops, expose configuration choices. \\
vCPU execution & \path{vcpu.c} & Owns KVM\_RUN, register marshaling, migration pinning, FPU capture, dispatch, and per-vCPU state. \\
Syscall path & \path{syscall_trap.c}, \path{lstar_gadget.S} & Handles LSTAR entry, in-guest fast syscalls, slow-path exits, and return-state repair. \\
Memory and regions & \path{memslot.c}, \path{region.c}, backend memory ops & Maps \UML{} memory changes into KVM-visible regions and memslot policy. \\
Exceptions & \path{exception.c} & Installs and handles IDT/TSS/IST paths, page faults, and special recovery paths. \\
State tracing & \path{state_trace.[ch]} & Gives low-overhead, static-key-controlled visibility into state ownership bugs. \\
Tests & \path{test_marshal.c}, \path{test_byteshape.c} & Protect register ABI, gadget byte shape, and internal assumptions. \\
\bottomrule
\end{tabularx}
\end{table}

\subsection{Design Detail: Fast Path Versus Slow Path}

The backend now has two different syscall cost classes. Trivial
syscalls covered by the LSTAR gadget stay inside the KVM guest and
return directly to guest userspace. Syscalls outside the gadget set,
or paths requiring host assistance, still exit to the host and run the
full backend machinery. The project should keep these paths described
separately:

\begin{itemize}
  \item \textbf{Fast path:} getpid-family and clock-like calls that
        can be answered from guest-visible state with preserved ABI
        registers and no host-side work.
  \item \textbf{Slow path:} mmap, munmap, brk, I/O, scheduling,
        signal, and fault paths that need backend state or host
        services.
  \item \textbf{Recovery path:} interrupted LSTAR, page-fault stub,
        \#NM, and signal delivery edges where the backend must rebuild
        user-visible architectural state exactly.
\end{itemize}

\subsection{State-Audit Method}

The state-audit directory is part of the design, not just project
diary. It converted a vague SMP corruption class into a finite audit:
149 state items, 50+ operations, an ownership matrix, suspect
register audits, runtime tracing, and ranked findings.

\begin{table}[h]
\centering
\caption{State-audit layers.}
\begin{tabularx}{\linewidth}{@{}p{0.09\linewidth}p{0.29\linewidth}X@{}}
\toprule
Layer & Artefact & Contribution \\
\midrule
1 & State inventory & Enumerates register, control, segment, MSR, FPU, event, descriptor, per-task, per-vCPU, and per-mm state. \\
2 & Operations inventory & Lists hardware, KVM, backend, scheduler, and signal operations that read or write state. \\
3 & Ownership matrix & States who owns each state item at each transition. \\
4 & Suspect register audits & Tracks RBP, RAX, RDX, FS\_BASE and similar corruption signatures. \\
5 & Toolkit & cscope, ripgrep, ast-grep, bpftrace, ftrace, and parser recipes. \\
6--8 & Findings and fixes & Converts candidates into patches and regression matrices. \\
\bottomrule
\end{tabularx}
\end{table}

The most important fix from this arc was the SMP-T13
\texttt{migrate\_disable()} finding. \UML{} was using
\texttt{preempt\_disable()} to pin access to per-host-CPU vCPU state,
but the build did not have preemption count semantics, so the call did
not actually prevent migration. Replacing that assumption with
\texttt{migrate\_disable()} removed a structural cross-vCPU ownership
violation.

\subsection{Recent Fixes}

\begin{table}[h]
\centering
\caption{Selected recent \KVMvTwo{} correctness and performance closures.}
\begin{tabularx}{\linewidth}{@{}p{0.17\linewidth}p{0.25\linewidth}X@{}}
\toprule
Item & Closure & Why it matters \\
\midrule
SMP-T41 & Recover user RAX in EINTR-mid-PF-stub & Closed the true mt-mini byte-zero residual. \\
SMP-T54 & SOCK\_CLOEXEC and \texttt{umlctl} sysrq halt & Removed worker fd leak and init shutdown hang class. \\
SMP-T55 & Per-vCPU FPU-dirty epoch & Restored perf-py-startup gate while preserving cross-task FPU guarantees. \\
SMP-T57 Phase A & CR4.OSXSAVE, XCR0=0x7, CPUID AVX/XSAVE unmask & Turned stress-ng AVX \#UD from a fragile mask problem into supported guest architectural state. \\
Gadget revival & LSTAR in-guest fast path & Collapsed trivial syscall cost from host-exit scale to roughly 90 cycles on measured hosts. \\
\bottomrule
\end{tabularx}
\end{table}

\subsection{Design Boundary}

The current backend should be presented upstream as a backend for
\UML{}, not as a general VMM. The VMM-like machinery exists only to
make \UML{} execution fast and faithful. That framing matters: it
keeps device emulation, boot firmware, hardware discovery, and cloud
microVM expectations out of scope.
